\documentclass[11pt,a4paper]{article}

% ---------------------------------------------------------------
% Packages
% ---------------------------------------------------------------
\usepackage[margin=2.5cm]{geometry}
\usepackage{amsmath,amssymb}
\usepackage{microtype}
\usepackage{parskip}
\usepackage{enumitem}
\usepackage{hyperref}
\usepackage{xcolor}
\usepackage{listings}
\usepackage{booktabs}

% ---------------------------------------------------------------
% Code listings style
% ---------------------------------------------------------------
\lstset{
  basicstyle=\ttfamily\small,
  breaklines=true,
  frame=single,
  keywordstyle=\color{blue},
  language=Python,
  commentstyle=\color{gray},
  showstringspaces=false,
}

% ---------------------------------------------------------------
% Logical notation macros
% ---------------------------------------------------------------
\newcommand{\Cn}{\mathrm{Cn}}
\newcommand{\entails}{\models}
\newcommand{\nentails}{\not\models}
\newcommand{\contraction}{\mathbin{\div}}
\newcommand{\revision}{\mathbin{*}}
\newcommand{\expansion}{\mathbin{+}}
\newcommand{\lneg}{\neg}

% ---------------------------------------------------------------
% Title
% ---------------------------------------------------------------
\title{\textbf{Belief Revision Engine}\\
       \large DTU 02180 – Introduction to Artificial Intelligence\\
       Group Assignment 2, Spring 2026}
\author{Group AI3}
\date{April 2026}

\begin{document}
\maketitle
\tableofcontents
\newpage

% ===============================================================
\section{Introduction}
% ===============================================================

A central challenge in knowledge representation is how to update an agent's
beliefs in the light of new information.
As posed in the course slides, \emph{``how do you update a database of knowledge
in the light of new information?''}
When an agent acquires a new observation $\varphi$ that contradicts some of its
existing beliefs, it cannot simply add $\varphi$ to its belief set without
first removing the conflicting material, because an inconsistent belief set
entails everything and is therefore useless.

We address this problem using the \emph{belief base approach}:
beliefs are represented as an ordered, finite set of propositional formulas
equipped with a \emph{priority ordering} that encodes degrees of entrenchment.
We implement three belief-change operations:

\begin{itemize}[noitemsep]
  \item \textbf{Expansion} $B \expansion \varphi$: add new information unconditionally.
  \item \textbf{Contraction} $B \contraction \varphi$: remove a belief and all formulas that
        jointly entail it, while preserving the most entrenched beliefs.
  \item \textbf{Revision} $B \revision \varphi$: incorporate potentially conflicting
        new information, defined via the \emph{Levi Identity}:
        $B \revision \varphi := (B \contraction \lneg\varphi) \expansion \varphi$.
\end{itemize}

Logical entailment is decided by the resolution algorithm of Russell \& Norvig
(2020, henceforth AIMA), implemented from scratch without any external SAT solver.
Contraction uses \emph{partial meet contraction} with a priority-based selection
function.  We verify the implementation against the five AGM revision postulates.

% ===============================================================
\section{Belief Base}
% ===============================================================

\subsection{Propositional language}

We work over the propositional language $\mathcal{L}$ generated by a countable
set of \emph{atoms} $\{p, q, r, \ldots\}$ and five connectives
$\lneg, \land, \lor, \to, \leftrightarrow$.
The set of well-formed formulas (WFF) is defined by the grammar
\[
  \varphi \;::=\; p \;\mid\; \lneg\varphi \;\mid\; (\varphi \land \varphi) \;\mid\;
  (\varphi \lor \varphi) \;\mid\; (\varphi \to \varphi) \;\mid\;
  (\varphi \leftrightarrow \varphi)
\]
following AIMA Figure 7.7.
Atoms are the atomic sentences; a \emph{literal} is an atom $p$ or its negation $\lneg p$.

\subsection{Belief base with priority ordering}

Following the belief-base approach (as opposed to the belief-set approach based on
plausibility orders, also discussed in the slides), we represent the agent's epistemic
state as a finite ordered list
\[
  B = [\varphi_0,\, \varphi_1,\, \ldots,\, \varphi_{n-1}]
\]
where the \emph{index} of a formula encodes its \emph{entrenchment}: $\varphi_0$ is
the most entrenched belief (it will be the last to be removed during contraction) and
$\varphi_{n-1}$ is the least entrenched.

\textbf{Example belief base} (from slides9, ``Bob''):
\[
  B = [p,\; q,\; (p \to q)]
  \quad \text{with priority } p > q > (p \to q).
\]
This represents the state where Bob believes $p$ with highest confidence,
$q$ with medium confidence, and $p \to q$ with the lowest confidence among the three.

% ===============================================================
\section{Logical Entailment via Resolution}
% ===============================================================

\subsection{Entailment}

We write $\mathit{KB} \entails \alpha$ to mean that $\alpha$ is a \emph{logical
consequence} of $\mathit{KB}$: in every model (truth-value assignment) in which
all formulas in $\mathit{KB}$ are true, $\alpha$ is also true (AIMA §7.3).
The logical closure of $B$ is $\Cn(B) = \{\varphi \mid B \entails \varphi\}$.

\subsection{CNF conversion}

The resolution procedure requires the input to be in
\emph{conjunctive normal form} (CNF): a conjunction of clauses, where each
clause is a disjunction of literals (AIMA §7.5.2).
We convert any formula to CNF in four steps (AIMA Figure 7.12):

\begin{enumerate}[noitemsep]
  \item \textbf{Eliminate biconditionals}: replace $(\alpha \leftrightarrow \beta)$
        with $(\alpha \to \beta) \land (\beta \to \alpha)$.
  \item \textbf{Eliminate implications}: replace $(\alpha \to \beta)$
        with $(\lneg\alpha \lor \beta)$.
  \item \textbf{Move negations inward}: apply double-negation elimination
        ($\lneg\lneg\alpha \equiv \alpha$) and De Morgan's laws:
        $\lneg(\alpha \land \beta) \equiv (\lneg\alpha \lor \lneg\beta)$,\quad
        $\lneg(\alpha \lor \beta) \equiv (\lneg\alpha \land \lneg\beta)$.
  \item \textbf{Distribute $\lor$ over $\land$}:
        $(\alpha \lor (\beta \land \gamma)) \equiv
         (\alpha \lor \beta) \land (\alpha \lor \gamma)$.
\end{enumerate}

\textbf{Worked example}: convert $B_{1,1} \leftrightarrow (P_{1,2} \lor P_{2,1})$
(from AIMA §7.5.2):
\begin{align*}
\text{Step 1:} &\quad (B_{1,1} \to (P_{1,2} \lor P_{2,1})) \land
                       ((P_{1,2} \lor P_{2,1}) \to B_{1,1}) \\
\text{Step 2:} &\quad (\lneg B_{1,1} \lor P_{1,2} \lor P_{2,1}) \land
                       (\lneg(P_{1,2} \lor P_{2,1}) \lor B_{1,1}) \\
\text{Step 3:} &\quad (\lneg B_{1,1} \lor P_{1,2} \lor P_{2,1}) \land
                       ((\lneg P_{1,2} \land \lneg P_{2,1}) \lor B_{1,1}) \\
\text{Step 4:} &\quad (\lneg B_{1,1} \lor P_{1,2} \lor P_{2,1}) \land
                       (\lneg P_{1,2} \lor B_{1,1}) \land (\lneg P_{2,1} \lor B_{1,1})
\end{align*}
The result is a conjunction of three clauses, each a disjunction of literals.

\subsection{Resolution rule and proof by refutation}

The \textbf{resolution rule} (AIMA §7.5.2) states: from a clause
$(\ell \lor C_1)$ and a clause $(\lneg\ell \lor C_2)$, we may derive the
\textbf{resolvent} $C_1 \lor C_2$, where $\ell$ and $\lneg\ell$ are
complementary literals.

To decide $\mathit{KB} \entails \alpha$, we use \textbf{proof by refutation}
(AIMA Fig. 7.13):
convert $(\mathit{KB} \land \lneg\alpha)$ to CNF,
then repeatedly apply the resolution rule to pairs of clauses,
adding each new resolvent to the clause set.
If the \emph{empty clause} (denoted $\square$, representing $\bot$) is derived,
then $\mathit{KB} \land \lneg\alpha$ is unsatisfiable, so $\mathit{KB} \entails \alpha$.
If no new clauses can be added (fixpoint reached), then $\mathit{KB} \nentails \alpha$.

The algorithm is \textbf{sound}: every resolvent is a logical consequence of
its parent clauses.
It is \textbf{complete} by the \emph{Ground Resolution Theorem} (AIMA §7.5.2):
if a set of clauses is unsatisfiable then its resolution closure contains the
empty clause.

\textbf{Small worked example}: check $\{p,\; p \to q\} \entails q$:
\begin{enumerate}[noitemsep]
  \item Add $\lneg q$ (negated query).  CNF clauses: $\{p\}$, $\{\lneg p, q\}$, $\{\lneg q\}$.
  \item Resolve $\{p\}$ with $\{\lneg p, q\}$: get $\{q\}$.
  \item Resolve $\{q\}$ with $\{\lneg q\}$: get $\square$ (empty clause).
  \item Empty clause derived $\Rightarrow$ $\{p, p\to q\} \entails q$. \quad\checkmark
\end{enumerate}

% ===============================================================
\section{Contraction}
% ===============================================================

\subsection{Remainder sets}

For a belief base $B$ and formula $\varphi$, the \textbf{remainder set}
$B \perp \varphi$ is the collection of all \emph{maximal} subsets of $B$
that do not entail $\varphi$:
\[
  B \perp \varphi \;=\; \bigl\{ S \subseteq B \;\mid\;
    S \nentails \varphi \;\text{ and }\; \forall\, S' \text{ s.t. } S \subsetneq S' \subseteq B:
    S' \entails \varphi \bigr\}.
\]
The \emph{maximality} condition ensures we remove as little as possible.

\subsection{Selection function and partial meet contraction}

We use a priority-based \textbf{selection function} $\gamma$: for each remainder
$S \in B \perp \varphi$, compute the sorted tuple of priority indices of its
members.  We select all remainders whose sorted index tuple is lexicographically
minimal (equivalently, that retain the most-entrenched beliefs).

\textbf{Partial meet contraction} is defined as:
\[
  B \contraction \varphi \;:=\; \bigcap \gamma(B \perp \varphi).
\]
That is, $B \contraction \varphi$ contains exactly the formulas present in
\emph{every} selected remainder.

\textbf{Worked example}: $B = [p,\; q,\; (p \to q)]$, contract by $q$.

Since $B \entails q$, contraction is non-vacuous.
The remainders $B \perp q$ are:
\begin{itemize}[noitemsep]
  \item $S_1 = \{p\}$\quad (removing $q$ and $(p\to q)$ yields a set not entailing $q$)
  \item $S_2 = \{(p \to q)\}$\quad (similarly)
  \item $S_3 = \{p, q \text{ removed, }(p\to q) \text{ removed}\} = \{p\}$ — same as $S_1$
\end{itemize}
More precisely, the maximal subsets of $B$ not entailing $q$ are:
$\{p\}$ and $\{(p \to q)\}$.
The selection function prefers $\{p\}$ (contains index-0 element $p$, the most entrenched).
Thus $\gamma(B \perp q) = \{\{p\}\}$ and $B \contraction q = \{p\}$.

\subsection{Contraction postulates (informal)}

Our implementation satisfies the following AGM contraction postulates:

\begin{description}[noitemsep,labelwidth=5em]
  \item[Failure] If $B \nentails \varphi$, then $B \contraction \varphi = B$ (vacuous
        contraction): implemented by an immediate check.
  \item[Success] $B \contraction \varphi \nentails \varphi$ (unless $\varphi$ is a
        tautology): guaranteed because we only keep subsets of $B$ that do not
        entail $\varphi$.
  \item[Inclusion] $B \contraction \varphi \subseteq B$: the contraction is a subset
        of the original base by construction.
  \item[Relevance] Every formula dropped in going from $B$ to $B \contraction \varphi$
        was relevant to entailing $\varphi$: guaranteed by the maximality of each
        remainder.
  \item[Uniformity] If any subset of $B$ that implies $\varphi$ also implies $\psi$,
        then $B \contraction \varphi$ and $B \contraction \psi$ keep the same formulas:
        satisfied since selection depends only on the structure of remainders, which is
        symmetric with respect to $\varphi$ and $\psi$ under this condition.
\end{description}

% ===============================================================
\section{Expansion}
% ===============================================================

Expansion is the simplest belief-change operation.  For a formula $\varphi$:
\[
  B \expansion \varphi \;:=\; \Cn(B \cup \{\varphi\}).
\]
In a finite syntactic implementation, we simply add $\varphi$ to the list of
formulas (at the lowest priority, i.e. appended at the end).
Expansion always succeeds in adding $\varphi$ to the belief base, but it
does not check for consistency: expanding with $\lneg q$ when $q \in B$
produces an inconsistent belief base.  This is why revision (which first
removes the obstacle) is the more appropriate operation when the incoming
information may conflict with existing beliefs.

% ===============================================================
\section{Revision}
% ===============================================================

\subsection{Levi Identity}

Revision is defined via the \textbf{Levi Identity} (slides9):
\[
  B \revision \varphi \;:=\; (B \contraction \lneg\varphi) \expansion \varphi.
\]
The rationale is to first remove from $B$ all formulas that would block the
acceptance of $\varphi$ (contraction by $\lneg\varphi$), and then add $\varphi$
(expansion).  This two-step process ensures that $\varphi$ is accepted
and that the result is consistent whenever $\varphi$ itself is consistent.

\subsection{AGM revision postulates}

We verify the five AGM revision postulates (slides9).
All 22 tests in our test suite pass:

\begin{description}[noitemsep]
  \item[(R1) Success] $\varphi \in \Cn(B \revision \varphi)$.
  \textit{Satisfied}: after revision, $\varphi$ is in the expanded base by
  definition of expansion.

  \item[(R2) Inclusion] $B \revision \varphi \subseteq B \expansion \varphi$.
  \textit{Satisfied}: the contracted base $(B \contraction \lneg\varphi)$ is
  a subset of $B$, so $(B \contraction \lneg\varphi) \expansion \varphi
  \subseteq B \expansion \varphi$.

  \item[(R3) Vacuity] If $\lneg\varphi \notin \Cn(B)$, then
  $B \revision \varphi = B \expansion \varphi$.
  \textit{Satisfied}: if $B \nentails \lneg\varphi$, then
  $B \contraction \lneg\varphi = B$ (Failure postulate of contraction),
  so $B \revision \varphi = B \expansion \varphi$.

  \item[(R4) Consistency] $B \revision \varphi$ is consistent unless $\varphi$
  is a contradiction.
  \textit{Satisfied}: after contracting $\lneg\varphi$, the base does not
  entail $\lneg\varphi$.  Adding $\varphi$ then yields a consistent base
  (since $\varphi$ is consistent and does not conflict with the contracted base).

  \item[(R5) Extensionality] If $\varphi \equiv \psi$ then
  $B \revision \varphi = B \revision \psi$.
  \textit{Satisfied}: logically equivalent formulas have the same negation
  (modulo equivalence), so $B \contraction \lneg\varphi$ and
  $B \contraction \lneg\psi$ yield equivalent bases, and adding equivalent
  $\varphi$ or $\psi$ preserves equivalence.
\end{description}

\subsection{Worked example}

Let $B = [p,\; q,\; (p \to q)]$ and revise by $\lneg q$:

\begin{enumerate}[noitemsep]
  \item \textbf{Contraction} $B \contraction q$ (since $\lneg(\lneg q) = q$):
        Remainders of $B \perp q$ are $\{p\}$ and $\{(p\to q)\}$.
        Selection function chooses $\{p\}$ (most entrenched).
        $B \contraction q = \{p\}$.
  \item \textbf{Expansion} $\{p\} \expansion \lneg q = \{p,\; \lneg q\}$.
  \item \textbf{Result}: $B \revision \lneg q = \{p,\; \lneg q\}$.
\end{enumerate}

Checking the postulates:
(R1) $\lneg q \in B \revision \lneg q$ \checkmark;\quad
(R4) $\{p, \lneg q\}$ is consistent \checkmark;\quad
(R3) not applicable since $B \entails q = \lneg(\lneg q)$ \checkmark.

% ===============================================================
\section{What We Learned}
% ===============================================================

Implementing the belief revision engine from scratch taught us several lessons.

\textbf{CNF and resolution.}
Implementing full CNF conversion correctly required careful attention to the
distribution of $\lor$ over $\land$.  Formulas such as
$(p \lor (q \land r)) \lor (s \land t)$ require multiple recursive passes.
We also had to handle tautological clauses (those containing both $p$ and
$\lneg p$) separately, since they would otherwise cause spurious entailments.

\textbf{Remainder-set enumeration.}
Finding all maximal non-entailing subsets is exponential in the size of the
belief base (there are up to $2^n$ subsets for a base of size $n$).
For small bases this is manageable, but for larger belief bases a more
efficient representation — e.g., a partial order over formulas — would be
necessary.  We also found it subtle to ensure \emph{maximality}: a naive
approach that removes one formula at a time does not always produce all
maximal subsets.

\textbf{Theory vs.\ implementation.}
The AGM framework is defined at the level of logical consequences ($\Cn(B)$),
while our implementation works at the syntactic level (finite lists of
formulas).  Some care is needed when comparing belief bases for semantic
equality (we check mutual entailment rather than syntactic identity).
The Levi Identity elegantly bridges contraction and revision, and its
correctness is easy to verify at the semantic level.  Translating this to
the syntactic level required showing that the priority-based selection
function tracks the intended semantic notion of ``most plausible world.''

% ===============================================================
\section{Conclusion and Further Work}
% ===============================================================

We have implemented a complete, working belief revision engine in Python using
only the standard library.  The engine supports:

\begin{itemize}[noitemsep]
  \item Full propositional formula parsing with all five connectives.
  \item CNF conversion following the four-step algorithm of AIMA Chapter 7.
  \item A resolution-based entailment checker (PL-Resolution, AIMA Fig. 7.13).
  \item Partial meet contraction with a priority-based selection function.
  \item Expansion and revision via the Levi Identity.
  \item An AGM postulate test suite (22 tests, all passing).
\end{itemize}

Several directions for further work suggest themselves:

\begin{description}[noitemsep]
  \item[Efficiency] The remainder-set enumeration is exponential.
        Clever indexing of clauses (which literals appear positive/negative
        in which clauses) and efficient data structures for the clause set
        would substantially speed up the resolution procedure.

  \item[Iterated revision] The current implementation performs single-step
        revision.  Iterated revision (applying revision repeatedly as new
        information arrives) requires a dynamic representation of entrenchment,
        e.g., an epistemic entrenchment ordering or a TPO (total pre-order over
        possible worlds) as discussed in the slides.

  \item[Plausibility-order approach] The slides present an alternative
        formulation of belief revision in terms of plausibility orders over
        valuations.  It would be interesting to implement this approach and
        compare it with the belief-base approach for the same examples.

  \item[Mastermind extension] The assignment mentions an optional Mastermind
        extension.  The belief base could model the game's constraints as
        propositional formulas, and revision could be used to narrow down
        possible solutions after each guess.
\end{description}

% ---------------------------------------------------------------
\section*{References}
% ---------------------------------------------------------------

\begin{enumerate}[label={[\arabic*]},noitemsep]
  \item Russell, S.\ and Norvig, P.\ (2020).
        \textit{Artificial Intelligence: A Modern Approach}, 4th ed.
        Chapter 7: Logical Agents.
  \item Gärdenfors, P.\ (1988).
        \textit{Knowledge in Flux: Modelling the Dynamics of Epistemic States}.
        MIT Press.  (As presented in DTU 02180 slides9.)
  \item DTU 02180 Lecture Slides, Slides 8 and 9:
        Propositional Logic, Resolution, Belief Revision (Spring 2026).
\end{enumerate}

\end{document}
